\chapter{Серия \texttt{7x}: перенос на эмбрион и локализация нового bottleneck}

\section{Почему перенос не был формальностью}

После \texttt{6d} легко было бы вообразить, что осталось просто применить
победивший гибрид к большому реальному графу. Но именно здесь и начинается самая
поучительная часть всей серии. Перенос показал, что формула, убедительная на
золотом мире, еще не гарантирует содержательного прогресса на эмбрионе. Причина
при этом оказалась не тривиальной.

\section{\texttt{exp7a}: end-to-end transfer check}

\subsection{Замысел}

\texttt{exp7a} спрашивал: сохраняются ли выводы \texttt{6c--6d} при переходе на
крупный seeded subgraph, где нет ручной capacity и где retrieval должен работать
на реальной, а не hand-authored топологии.

Сравнивались raw diffusion, in-strength prior, PageRank prior, local leak only и
hybrid local leak + PageRank.

\subsection{Стенд}

Подграф строился как seeded expansion по понятиям, продолжающим мотивы
\texttt{6x}: кофе, паспорт, apartment, charger, travel, morning и т.\,д.
Формально ceiling на concept nodes составлял 10\,000, но включение ATOMIC
создало гораздо более массивную оболочку episode nodes.

\subsection{Что получилось вместо чистого transfer-benchmark}

Фактический размер графа оказался таким:

\begin{itemize}[leftmargin=2em]
\item 10\,000 concept nodes;
\item 455\,481 episode nodes;
\item 465\,481 узел суммарно.
\end{itemize}

Уже одно это сильно меняет характер задачи. Но еще важнее был другой эффект:
automatic anchoring порождал шумные source clouds вроде \texttt{thing},
\texttt{matter}, \texttt{before}, \texttt{early}, \texttt{imbibement},
\texttt{morningness}.

\subsection{Результат}

Все tested methods дали одинаковую дискретную картину:

\begin{itemize}[leftmargin=2em]
\item 4/13 pairwise checks;
\item 3/7 solved queries.
\end{itemize}

Численные margins различались сильно, но ordering не менялся.

\subsection{Почему это не был ``нулевой'' опыт}

На первый взгляд такой tie можно было бы счесть провалом. Но на деле
\texttt{exp7a} дал один из важнейших диагнозов серии:

\begin{quote}
в реальном end-to-end режиме retrieval law почти не тестируется, потому что
pipeline уже ломается на этапе query $\to$ source injection.
\end{quote}

Кроме того, top-k были засорены echo самих source nodes, так что наблюдаемое
``retrieval'' часто было просто возвратом к уже впрыснутым якорям.

\begin{takeaway}
\texttt{exp7a} не опроверг local+global hybrid. Он показал, что на реальном
эмбрионе более ранний слой -- anchoring и source injection -- доминирует над
всем, что происходит позже. Это и стало основанием для controlled-source опыта
\texttt{7b}.
\end{takeaway}

\section{\texttt{exp7b}: controlled-source transfer}

\subsection{Почему это новый эксперимент, а не ``починка 7a''}

\texttt{7b} делает стандартный физический ход: замораживает вход, чтобы
проверять среду. Manual anchors здесь не подаются как финальная архитектура для
natural-language retrieval. Они используются как калиброванный впрыск сигнала.

Именно поэтому \texttt{7a} и \texttt{7b} должны читаться вместе:

\begin{itemize}[leftmargin=2em]
\item \texttt{7a} локализует bottleneck automatic anchoring;
\item \texttt{7b} временно убирает этот bottleneck из игры.
\end{itemize}

\subsection{Два варианта графа}

Отдельно сравнивались:

\begin{itemize}[leftmargin=2em]
\item ConceptNet only;
\item ConceptNet + ATOMIC.
\end{itemize}

Это было методически важно: после \texttt{7a} нельзя было молча допустить, что
гигантская ATOMIC shell никак не меняет topology priors.

\subsection{\texttt{7b.1}: первая controlled-source версия}

В первой версии manual cues были чистыми, но во многом еще scene-like:
\texttt{check in + airport}, \texttt{home + remember}, \texttt{morning +
breakfast}.

Результат: снова tie across methods.

\begin{itemize}[leftmargin=2em]
\item pairwise: 2/12;
\item solved: 2/7.
\end{itemize}

Главный урок этого захода: даже ручные anchors бесполезны, если сами cues
слишком общие и описывают сцену, а не механизм ответа.

\subsection{\texttt{7b.2}: diagnostic-cue version}

Во второй версии cues сделали более механизмо-ориентированными:

\begin{itemize}[leftmargin=2em]
\item \texttt{coffee + stimulant} вместо \texttt{coffee + energy};
\item \texttt{check in counter + document} вместо \texttt{check in + airport};
\item \texttt{wake up + tired} вместо \texttt{morning + breakfast}.
\end{itemize}

Это улучшило семантическую осмысленность top-k, но discrete separation все равно
не появилось.

\subsection{Результат \texttt{7b.2}}

На обеих graph variants получились одинаковые агрегаты:

\begin{itemize}[leftmargin=2em]
\item primary pairwise: 1/8;
\item primary solved: 1/5;
\item control pairwise: 2/4;
\item control solved: 1/2.
\end{itemize}

То есть controlled-source regime оказался методически правильным, но формулы
снова не разошлись.

\subsection{Что именно теперь стало ясно}

После \texttt{7b} уже нельзя было объяснять стагнацию одной лишь проблемой
automatic anchors. Bottleneck сместился глубже:

\begin{enumerate}[leftmargin=2em]
\item либо signal идет в broad intermediate hubs;
\item либо он попадает в graph-native adjacent manifold, а не в hand-named
target node;
\item либо оба эффекта происходят одновременно.
\end{enumerate}

Примеры этого были вполне содержательны:

\begin{itemize}[leftmargin=2em]
\item для morning query граф охотно выдавал \texttt{drink coffee}, а не обязательно
\texttt{coffee} или \texttt{espresso};
\item для laptop query граф выдавал \texttt{power}, \texttt{battery},
\texttt{generator}-family вместо \texttt{charger};
\item для check-in query граф поднимал procedural nodes вроде
\texttt{checking in} сильнее, чем \texttt{passport}.
\end{itemize}

\begin{riskbox}
Здесь очень легко сделать ложный вывод ``PageRank бесполезен'' или
``local leak не переносится''. На самом деле \texttt{7b} показал более узкую
вещь: оба механизма работают слишком поздно в pipeline, если transport already
lands in the wrong manifold.
\end{riskbox}

\section{\texttt{exp7c}: conjunctive cue-overlap}

\subsection{Новый вопрос}

После повторяющихся tie в \texttt{7b} естественно возникла новая гипотеза:
может быть, проблема не в знаменателе вообще, а в том, как складываются cue
fields. Если broad hub ярко светится от одного сильного cue, а правильный ответ
получает умеренную поддержку от нескольких независимых cues, то простой сумматор
может систематически переоценивать первый случай.

Отсюда и возникает overlap idea:

\[
\mathrm{Overlap}_{\mathrm{geom}}(v)=
\left(\prod_g (A_g(v)+\varepsilon)\right)^{1/G}.
\]

\subsection{Что проверялось}

Сравнивались:

\begin{itemize}[leftmargin=2em]
\item sum aggregation;
\item sum + PageRank;
\item geometric overlap;
\item geometric overlap + PageRank.
\end{itemize}

То есть это уже не еще один denominator benchmark, а новая попытка изменить
саму логику совместного учета нескольких cue-conditioned полей.

\subsection{Результат}

Именно здесь серия впервые довольно жестко исчерпала надежду на ``еще одну
формулу финального score''.

По дискретным метрикам sum-based и overlap-based варианты снова tied. В
исполненном \texttt{7c} это было зафиксировано как одинаковые primary outcomes
на обеих graph variants; в дальнейшей постановке \texttt{8a} этот итог
формулируется как устойчивый tie по уровням порядка \(4/10\) pairwise и \(3/6\)
solved для главных конкурентов.

\subsection{Почему это содержательно важно}

Overlap failure означает следующее:

\begin{quote}
нельзя ``перемножить'' путь, которого вообще нет. Если $A_g(v)=0$ для нужного
cue на нужной цели, то никакая интерференция или coincidence scoring не создаст
там содержательного сигнала.
\end{quote}

Это очень сильная отрицательная находка. Она переводит разговор от финального
скоринга к самой транспортной архитектуре.

\section{Что серия \texttt{7x} установила}

\subsection{Положительные результаты}

\begin{enumerate}[leftmargin=2em]
\item Automatic anchoring действительно является самостоятельным bottleneck.
\item Controlled-source regime нужен и методически оправдан.
\item Эффект source echo в retrieval-списках нужно явно устранять.
\item ATOMIC shell сильно влияет на размеры и priors, но не обязана менять сами
ordering outcomes.
\end{enumerate}

\subsection{Отрицательные результаты}

\begin{enumerate}[leftmargin=2em]
\item На реальном эмбрионе denominator families не дают такого separation, как на
золотых мирах.
\item Простая cue-overlap aggregation тоже не решает проблему.
\end{enumerate}

\subsection{Главный диагностический вывод}

К концу \texttt{7x} становится правдоподобной новая общая картина:

\begin{quote}
узкое место находится не столько в пост-фактум нормировке уже пришедшего
сигнала, сколько в том, что сигнал либо не доходит до правильного manifold,
либо доходит не теми путями, либо приходит в слишком соседний, а не точный
ответный узел.
\end{quote}

Отсюда уже естественным образом рождается серия \texttt{8a--9a}, где вопрос
смещается с denominators на transport wall.
